\section{Shipping}
\label{sec:shipping}

\subsection{Scope and inventory mapping}

This block targets the area-source component of CAMS GNFR~G
(\emph{Shipping}). In the CAMS-REG-ANT inventory, shipping emissions are
reported largely over marine areas, with most support concentrated in a
relatively small number of offshore parent cells over the Mediterranean,
Atlantic, and adjacent seas. The methodological problem at this stage is
therefore not detailed local downscaling, but the construction of a
fine-grid \emph{within-cell spatial weighting raster} that can later be
combined with the corresponding CAMS cell totals.

This distinction is essential. The workflow described here produces a
single-band raster of normalised fine-grid weights such that the
fine-pixel shares sum to unity within each CAMS GNFR~G area cell. The
layer does not by itself represent emissions. Its role is purely
allocative: to distribute each CAMS shipping cell total across plausible
shipping lanes and port environments on the common fine grid.

\subsection{Conceptual design}

The shipping proxy is intentionally simple and robust. Unlike sectors
whose internal processes require pollutant-specific or subsector-specific
differentiation, GNFR~G is treated here with one common spatial logic for
all pollutants. This choice reflects both the limited number of CAMS
shipping cells over the study domain and the fact that the main spatial
question is not process disaggregation, but the recovery of a realistic
marine--coastal support pattern inside those coarse cells.

The design therefore combines three complementary ingredients:
\begin{itemize}
  \item vessel-density information to represent the main offshore traffic
        corridors;
  \item port-related land-cover support to anchor part of the signal at
        the land--sea interface;
  \item mapped port infrastructure geometry to refine the coastal anchor.
\end{itemize}

This produces a transparent first-order weighting system. Maritime
traffic provides the dominant spatial signal, while port layers stabilise
the proxy near coastlines and major harbour environments. The result is a
proxy that remains consistent with the marine character of shipping
emissions without attempting an artificial level of detail not supported
by the parent inventory.

\subsection{Input datasets}

The main inputs are summarised in
Table~\ref{tab:shipping_inputs}. Each plays a distinct role in the
construction of the weighting raster.

\begin{table}[htbp]
  \centering
  \caption{Main inputs for CAMS GNFR~G shipping weighting layers.}
  \label{tab:shipping_inputs}
  \begin{tabular}{p{3.6cm}p{9.1cm}}
    \hline
    Ingredient & Role \\
    \hline
    CAMS-REG-ANT NetCDF & Parent-cell geometry for GNFR~G area sources and support for within-cell normalisation \\
    EMODnet vessel density & Main offshore traffic-intensity signal used to recover shipping-lane structure \\
    CORINE CLC 2018 & Common fine grid and land-cover support; class~123 identifies port areas \\
    OpenStreetMap port infrastructure & Geometric refinement of harbour and port-support areas at the coast \\
    NUTS geometry & Construction of the common proxy-grid domain \\
    \hline
  \end{tabular}
\end{table}

The CAMS grid defines the parent support within which weights must sum to
unity. EMODnet provides the most direct harmonised indicator of maritime
traffic intensity and therefore acts as the main driver of the proxy.
CORINE contributes a stable land-cover representation of port areas,
while OpenStreetMap adds more detailed coastal geometry where explicit
port structures are mapped.

\subsection{Proxy construction}

Let $i$ denote a fine pixel on the common reference grid. The method
constructs one non-negative raw proxy surface $P(i)$ from harmonised
marine and port-support layers. This surface does not represent
emissions directly; it represents the relative spatial plausibility of
shipping activity within each CAMS parent cell.

\subsubsection{EMODnet vessel density}

The EMODnet annual-average vessel-density raster is used as the primary
spatial driver. After reprojection and resampling to the common fine
grid, the corresponding shipping-density field is written as $D(i)$. A
normalised marine-traffic term is then defined by
\begin{equation}
  D_n(i) = \frac{D(i)}{\max_k D(k)}.
\end{equation}

This term captures the broad structure of maritime traffic, including
major shipping lanes and frequently used marine corridors. Because the
target of the workflow is a weighting raster rather than a direct traffic
model, the EMODnet layer is interpreted only as a relative intensity
field. In the implementation, its influence is damped over clearly
terrestrial pixels so that offshore density values do not dominate
well inland.

\subsubsection{Port areas from CORINE}

CORINE Land Cover provides the common reference grid used throughout the
proxy workflow and supplies a stable land-based anchor through class~123
(\emph{port areas}). Let $CLC(i)$ denote the fractional support of port
areas in pixel $i$ after resampling to the fine grid. This term is not
intended to describe ship traffic intensity. Its role is to identify
locations where shipping-related emissions can plausibly interact with
coastal and harbour environments.

\subsubsection{OpenStreetMap port infrastructure}

OpenStreetMap complements CORINE with more detailed coastal geometry. The
retained features correspond to explicit port and harbour environments,
namely polygons or linear structures tagged as \texttt{landuse=port},
\texttt{harbour=yes}, \texttt{natural=harbour}, and
\texttt{man\_made=quay}, \texttt{man\_made=pier}, or
\texttt{man\_made=breakwater}. Candidate industrial tags such as
\texttt{industrial=shipyard} were considered at screening stage but are
not retained in the final proxy in order to avoid extending the shipping
signal onto generic industrial land that is not necessarily part of the
port interface.

Let $OSM(i)$ denote the fractional area of selected OpenStreetMap port
features within fine pixel $i$:
\begin{equation}
  OSM(i) =
  \frac{
    \mathrm{area}\!\left(
      \text{selected OSM port features} \cap i
    \right)
  }{
    \mathrm{area}(i)
  }.
\end{equation}

This term acts as a detailed harbour-support layer rather than as a
standalone estimate of activity intensity.

\subsubsection{Combined proxy}

The three ingredients are converted to a common unitless scale before
combination. For the OSM and CORINE port-support layers, a min--max
normalisation $z(\cdot)$ is applied on the fine grid. The raw combined
shipping proxy is then
\begin{equation}
  P(i) =
  0.25 \cdot D_n(i)
  + 0.5 \cdot z\!\left(OSM(i)\right)
  + 0.25 \cdot z\!\left(CLC(i)\right).
\end{equation}

This formulation assigns distinct roles to the three ingredients.
EMODnet preserves the dominant offshore traffic structure, while the OSM
and CORINE terms anchor the proxy to realistic coastal and port
locations. The weighting is intentionally conservative: it enriches the
marine signal with harbour support, but does not attempt to infer
subsector-specific behaviour or port-throughput intensity from the
ancillary datasets.

\subsection{Special considerations}

Several methodological choices follow directly from the character of the
shipping inventory. First, the number of relevant CAMS parent cells is
small compared with strongly land-based sectors, so a highly complex
downscaling design would add little practical value. Second, coastal
consistency is important: a purely marine proxy would underrepresent
ports, while a purely land-based proxy would misplace emissions that are
mainly generated offshore. Third, the proxy is designed to avoid
overfitting. It uses broad, harmonised datasets with clear physical
interpretation rather than attempting to reconstruct detailed operational
shipping behaviour unsupported by the parent inventory.

\subsection{Within-cell normalisation}

Let $m$ denote a CAMS GNFR~G area cell and let $i \in m$ denote fine
pixels whose centres fall inside that parent cell. The final within-cell
weight is
\begin{equation}
  W_G(i) =
  \frac{P(i)}
       {\displaystyle\sum_{k \in m} P(k)}.
\end{equation}

By construction,
\begin{equation}
  \sum_{i \in m} W_G(i) = 1
  \qquad \forall m.
\end{equation}

This normalisation ensures strict consistency with CAMS emission totals:
the weighting raster determines only how each parent-cell total is
distributed internally on the fine grid, without modifying the total mass
reported by the inventory.

\subsection{Fallback logic}

If the raw proxy support is absent in a given CAMS parent cell, that is,
if
\begin{equation}
  \sum_{k \in m} P(k) = 0,
\end{equation}
the method falls back to a uniform distribution over the valid fine-grid
support of that CAMS cell. This minimal fallback is sufficient for the
shipping case because the proxy is used only to preserve a complete
within-cell weighting field where ancillary evidence is locally absent.

\subsection{Output and interpretation}

The final product is a single-layer fine-grid raster,
\texttt{Shipping\_areasource.tif}, aligned with the common proxy grid used
for the other sectors. Each pixel stores a unitless within-cell weight
only. The raster should therefore be interpreted strictly as a spatial
allocation layer: it indicates how the CAMS GNFR~G area-source total
should be distributed inside each parent CAMS cell, but it does not by
itself represent downscaled emissions.
